\chapter{Blocks}\label{record-018}

\noindent{\footnotesize\color{BookMuted}\path{docs/design-system/blocks.md}\par}\smallskip

This document defines what Bijou means by a reusable block.

Blocks are larger, opinionated assemblies built from canonical patterns and component families.

\section{What a block is}\label{record-018-what-a-block-is}

A block is not:

\begin{itemize}
\tightlist
\item
  a one-off demo composition
\item
  a single primitive component
\item
  a shell-specific hack that only one proving surface understands
\end{itemize}

A block is:

\begin{itemize}
\tightlist
\item
  a reusable product-level assembly
\item
  composed from canonical patterns
\item
  built from existing component families and TUI ownership rules
\item
  documented with clear jobs, ownership, and lowering expectations
\end{itemize}

Blocks come after doctrine, patterns, and components.

That ordering matters:

\begin{enumerate}
\def\labelenumi{\arabic{enumi}.}
\tightlist
\item
  doctrine explains what good Bijou UX should feel like
\item
  patterns explain recurring interaction language
\item
  component families provide reusable primitives
\item
  blocks compose those primitives into repeatable product structures
\end{enumerate}

If a block appears before the first three layers are clear, it usually hardcodes accidental taste instead of reusable language.

\section{What blocks should standardize}\label{record-018-what-blocks-should-standardize}

Blocks should standardize:

\begin{itemize}
\tightlist
\item
  composition
\item
  structure
\item
  copy rhythm
\item
  focus ownership
\item
  common shell or workflow affordances
\end{itemize}

Blocks should not hide:

\begin{itemize}
\tightlist
\item
  important interaction ownership
\item
  lowering behavior
\item
  accessibility semantics
\item
  localization or directional assumptions
\end{itemize}

\section{Current public contract}\label{record-018-current-public-contract}

Bijou now exposes a metadata-first block authoring contract from \texttt{@\allowbreak{}f\allowbreak{}l\allowbreak{}y\allowbreak{}i\allowbreak{}n\allowbreak{}g\allowbreak{}r\allowbreak{}o\allowbreak{}b\allowbreak{}o\allowbreak{}t\allowbreak{}s\allowbreak{}/\allowbreak{}b\allowbreak{}i\allowbreak{}j\allowbreak{}o\allowbreak{}u}:

\begin{itemize}
\tightlist
\item
  \texttt{B\allowbreak{}l\allowbreak{}o\allowbreak{}c\allowbreak{}k\allowbreak{}M\allowbreak{}e\allowbreak{}t\allowbreak{}a\allowbreak{}d\allowbreak{}a\allowbreak{}t\allowbreak{}a} declares a block's package, name, family, scale, modes, slots, variants, config options, composed components, semantic facts, story ids, and docs.
\item
  \texttt{B\allowbreak{}l\allowbreak{}o\allowbreak{}c\allowbreak{}k\allowbreak{}D\allowbreak{}e\allowbreak{}f\allowbreak{}i\allowbreak{}n\allowbreak{}i\allowbreak{}t\allowbreak{}i\allowbreak{}o\allowbreak{}n} pairs that metadata with an explicit render function.
\item
  \texttt{B\allowbreak{}l\allowbreak{}o\allowbreak{}c\allowbreak{}k\allowbreak{}P\allowbreak{}a\allowbreak{}c\allowbreak{}k\allowbreak{}a\allowbreak{}g\allowbreak{}e\allowbreak{}M\allowbreak{}a\allowbreak{}n\allowbreak{}i\allowbreak{}f\allowbreak{}e\allowbreak{}s\allowbreak{}t} declares the exported blocks, package version, Bijou peer range, docs, and tags for ordinary NPM distribution.
\item
  \texttt{d\allowbreak{}e\allowbreak{}f\allowbreak{}i\allowbreak{}n\allowbreak{}e\allowbreak{}B\allowbreak{}l\allowbreak{}o\allowbreak{}c\allowbreak{}k\allowbreak{}(\allowbreak{})} and \texttt{d\allowbreak{}e\allowbreak{}f\allowbreak{}i\allowbreak{}n\allowbreak{}e\allowbreak{}B\allowbreak{}l\allowbreak{}o\allowbreak{}c\allowbreak{}k\allowbreak{}P\allowbreak{}a\allowbreak{}c\allowbreak{}k\allowbreak{}a\allowbreak{}g\allowbreak{}e\allowbreak{}(\allowbreak{})} validate definitions before authors publish or index them.
\item
  \texttt{d\allowbreak{}e\allowbreak{}f\allowbreak{}i\allowbreak{}n\allowbreak{}e\allowbreak{}B\allowbreak{}l\allowbreak{}o\allowbreak{}c\allowbreak{}k\allowbreak{}(\allowbreak{})} can attach DX-034 view data contracts and command intents for discovery without provider handles or command callbacks.
\item
  \texttt{a\allowbreak{}p\allowbreak{}p\allowbreak{}S\allowbreak{}h\allowbreak{}e\allowbreak{}l\allowbreak{}l\allowbreak{}B\allowbreak{}l\allowbreak{}o\allowbreak{}c\allowbreak{}k}, \texttt{r\allowbreak{}e\allowbreak{}a\allowbreak{}d\allowbreak{}e\allowbreak{}r\allowbreak{}S\allowbreak{}u\allowbreak{}r\allowbreak{}f\allowbreak{}a\allowbreak{}c\allowbreak{}e\allowbreak{}B\allowbreak{}l\allowbreak{}o\allowbreak{}c\allowbreak{}k}, \texttt{i\allowbreak{}n\allowbreak{}s\allowbreak{}p\allowbreak{}e\allowbreak{}c\allowbreak{}t\allowbreak{}o\allowbreak{}r\allowbreak{}P\allowbreak{}a\allowbreak{}n\allowbreak{}e\allowbreak{}l\allowbreak{}B\allowbreak{}l\allowbreak{}o\allowbreak{}c\allowbreak{}k}, \texttt{i\allowbreak{}n\allowbreak{}l\allowbreak{}i\allowbreak{}n\allowbreak{}e\allowbreak{}S\allowbreak{}t\allowbreak{}a\allowbreak{}t\allowbreak{}u\allowbreak{}s\allowbreak{}B\allowbreak{}l\allowbreak{}o\allowbreak{}c\allowbreak{}k}, \texttt{i\allowbreak{}n\allowbreak{}F\allowbreak{}l\allowbreak{}o\allowbreak{}w\allowbreak{}S\allowbreak{}t\allowbreak{}a\allowbreak{}t\allowbreak{}u\allowbreak{}s\allowbreak{}B\allowbreak{}l\allowbreak{}o\allowbreak{}c\allowbreak{}k}, \texttt{t\allowbreak{}r\allowbreak{}a\allowbreak{}n\allowbreak{}s\allowbreak{}i\allowbreak{}e\allowbreak{}n\allowbreak{}t\allowbreak{}O\allowbreak{}v\allowbreak{}e\allowbreak{}r\allowbreak{}l\allowbreak{}a\allowbreak{}y\allowbreak{}B\allowbreak{}l\allowbreak{}o\allowbreak{}c\allowbreak{}k}, \texttt{a\allowbreak{}c\allowbreak{}t\allowbreak{}i\allowbreak{}v\allowbreak{}i\allowbreak{}t\allowbreak{}y\allowbreak{}S\allowbreak{}t\allowbreak{}r\allowbreak{}e\allowbreak{}a\allowbreak{}m\allowbreak{}B\allowbreak{}l\allowbreak{}o\allowbreak{}c\allowbreak{}k}, \texttt{s\allowbreak{}h\allowbreak{}o\allowbreak{}r\allowbreak{}t\allowbreak{}c\allowbreak{}u\allowbreak{}t\allowbreak{}C\allowbreak{}u\allowbreak{}e\allowbreak{}B\allowbreak{}l\allowbreak{}o\allowbreak{}c\allowbreak{}k}, \texttt{p\allowbreak{}r\allowbreak{}o\allowbreak{}g\allowbreak{}r\allowbreak{}e\allowbreak{}s\allowbreak{}s\allowbreak{}I\allowbreak{}n\allowbreak{}d\allowbreak{}i\allowbreak{}c\allowbreak{}a\allowbreak{}t\allowbreak{}o\allowbreak{}r\allowbreak{}B\allowbreak{}l\allowbreak{}o\allowbreak{}c\allowbreak{}k}, the DOGFOOD component-family block exports, \texttt{s\allowbreak{}t\allowbreak{}a\allowbreak{}n\allowbreak{}d\allowbreak{}a\allowbreak{}r\allowbreak{}d\allowbreak{}B\allowbreak{}l\allowbreak{}o\allowbreak{}c\allowbreak{}k\allowbreak{}s}, \texttt{s\allowbreak{}t\allowbreak{}a\allowbreak{}n\allowbreak{}d\allowbreak{}a\allowbreak{}r\allowbreak{}d\allowbreak{}B\allowbreak{}l\allowbreak{}o\allowbreak{}c\allowbreak{}k\allowbreak{}S\allowbreak{}t\allowbreak{}o\allowbreak{}r\allowbreak{}i\allowbreak{}e\allowbreak{}s}, and \texttt{s\allowbreak{}t\allowbreak{}a\allowbreak{}n\allowbreak{}d\allowbreak{}a\allowbreak{}r\allowbreak{}d\allowbreak{}B\allowbreak{}l\allowbreak{}o\allowbreak{}c\allowbreak{}k\allowbreak{}P\allowbreak{}a\allowbreak{}c\allowbreak{}k\allowbreak{}a\allowbreak{}g\allowbreak{}e\allowbreak{}M\allowbreak{}a\allowbreak{}n\allowbreak{}i\allowbreak{}f\allowbreak{}e\allowbreak{}s\allowbreak{}t} expose the first-party block definitions: AppShell, ReaderSurface, InspectorPanel, inline status, in-flow status, transient overlay feedback, activity stream, shortcut cues, progress, grouping, explainability, documents, links, dividers, text entry, choices, navigation, disclosure, path progress, brand emphasis, mode-aware primitives, dense comparisons, hierarchies, exploration lists, and temporal dependency views. Their render functions now produce deterministic first-proof output across interactive, static, pipe, and accessible modes while keeping provider lifecycle, command dispatch, and production AppShell behavior out of the block definition layer.
\item
  \texttt{d\allowbreak{}e\allowbreak{}f\allowbreak{}i\allowbreak{}n\allowbreak{}e\allowbreak{}A\allowbreak{}p\allowbreak{}p\allowbreak{}S\allowbreak{}h\allowbreak{}e\allowbreak{}l\allowbreak{}l\allowbreak{}C\allowbreak{}o\allowbreak{}m\allowbreak{}p\allowbreak{}o\allowbreak{}s\allowbreak{}i\allowbreak{}t\allowbreak{}i\allowbreak{}o\allowbreak{}n\allowbreak{}(\allowbreak{})} groups runtime-backed blocks into logical navigation, content, inspector, status, and overlay slots without rendering.
\item
  \texttt{b\allowbreak{}l\allowbreak{}o\allowbreak{}c\allowbreak{}k\allowbreak{}R\allowbreak{}e\allowbreak{}n\allowbreak{}d\allowbreak{}e\allowbreak{}r\allowbreak{}N\allowbreak{}o\allowbreak{}d\allowbreak{}e\allowbreak{}(\allowbreak{})} and \texttt{r\allowbreak{}e\allowbreak{}n\allowbreak{}d\allowbreak{}e\allowbreak{}r\allowbreak{}B\allowbreak{}l\allowbreak{}o\allowbreak{}c\allowbreak{}k\allowbreak{}T\allowbreak{}r\allowbreak{}e\allowbreak{}e\allowbreak{}(\allowbreak{})} explicitly render nested block declarations when a host or preview wants child block output inside parent slots. Ordinary \texttt{b\allowbreak{}l\allowbreak{}o\allowbreak{}c\allowbreak{}k\allowbreak{}.\allowbreak{}r\allowbreak{}e\allowbreak{}n\allowbreak{}d\allowbreak{}e\allowbreak{}r\allowbreak{}(\allowbreak{})} remains a local single-block render contract.
\item
  \texttt{d\allowbreak{}e\allowbreak{}f\allowbreak{}i\allowbreak{}n\allowbreak{}e\allowbreak{}B\allowbreak{}l\allowbreak{}o\allowbreak{}c\allowbreak{}k\allowbreak{}S\allowbreak{}c\allowbreak{}h\allowbreak{}e\allowbreak{}m\allowbreak{}a\allowbreak{}A\allowbreak{}d\allowbreak{}a\allowbreak{}p\allowbreak{}t\allowbreak{}e\allowbreak{}r\allowbreak{}(\allowbreak{})}, \texttt{d\allowbreak{}e\allowbreak{}f\allowbreak{}i\allowbreak{}n\allowbreak{}e\allowbreak{}S\allowbreak{}c\allowbreak{}h\allowbreak{}e\allowbreak{}m\allowbreak{}a\allowbreak{}B\allowbreak{}l\allowbreak{}o\allowbreak{}c\allowbreak{}k\allowbreak{}(\allowbreak{})}, and \texttt{b\allowbreak{}i\allowbreak{}n\allowbreak{}d\allowbreak{}S\allowbreak{}c\allowbreak{}h\allowbreak{}e\allowbreak{}m\allowbreak{}a\allowbreak{}B\allowbreak{}l\allowbreak{}o\allowbreak{}c\allowbreak{}k\allowbreak{}I\allowbreak{}n\allowbreak{}p\allowbreak{}u\allowbreak{}t\allowbreak{}(\allowbreak{})} validate unknown boundary data before producing block render input and binding facts.
\item
  \texttt{v\allowbreak{}a\allowbreak{}l\allowbreak{}i\allowbreak{}d\allowbreak{}a\allowbreak{}t\allowbreak{}e\allowbreak{}B\allowbreak{}l\allowbreak{}o\allowbreak{}c\allowbreak{}k\allowbreak{}M\allowbreak{}e\allowbreak{}t\allowbreak{}a\allowbreak{}d\allowbreak{}a\allowbreak{}t\allowbreak{}a\allowbreak{}(\allowbreak{})}, \texttt{b\allowbreak{}l\allowbreak{}o\allowbreak{}c\allowbreak{}k\allowbreak{}M\allowbreak{}e\allowbreak{}t\allowbreak{}a\allowbreak{}d\allowbreak{}a\allowbreak{}t\allowbreak{}a\allowbreak{}S\allowbreak{}u\allowbreak{}m\allowbreak{}m\allowbreak{}a\allowbreak{}r\allowbreak{}y\allowbreak{}(\allowbreak{})}, and report helpers give tests, docs, DOGFOOD, MCP payloads, and agents a deterministic way to inspect block facts without rendering a terminal surface.
\end{itemize}

Blocks are described before they render, package authors avoid hidden global registration, and tooling can discover slots and mode support directly. Schema-bound blocks validate unknown boundary data without making Zod, GraphQL, or database schemas part of the core block dependency surface. Provider subscriptions, active hierarchy traversal, rendered AppShell, and DOGFOOD block multi-mode captures remain follow-on work.

\section{First credible block candidates}\label{record-018-first-credible-block-candidates}

\subsection{App frame}\label{record-018-app-frame}

The framed shell is already one of Bijou's strongest recurring structures.

As a future block, \texttt{A\allowbreak{}p\allowbreak{}p\allowbreak{} \allowbreak{}f\allowbreak{}r\allowbreak{}a\allowbreak{}m\allowbreak{}e} should standardize:

\begin{itemize}
\tightlist
\item
  pane rhythm
\item
  active-region treatment
\item
  shell header/footer ownership
\item
  shell overlays such as settings and help
\end{itemize}

\subsection{Settings drawer}\label{record-018-settings-drawer}

The settings drawer is a credible future block because it already has:

\begin{itemize}
\tightlist
\item
  canonical preference rows
\item
  shell-owned dismissal and discovery behavior
\item
  stable section rhythm
\item
  truthful active-row treatment
\end{itemize}

As a block, it should be reusable without depending on DOGFOOD-specific code.

\subsection{Notification center}\label{record-018-notification-center}

The notification center is a likely future block once Humane Terminal work resumes.

It should eventually standardize:

\begin{itemize}
\tightlist
\item
  notice rows
\item
  archive/history review
\item
  filter rhythm
\item
  shell-owned interrupt versus review behavior
\end{itemize}

\subsection{Guided flow}\label{record-018-guided-flow}

Guided flow is the general block family for:

\begin{itemize}
\tightlist
\item
  onboarding
\item
  setup sequences
\item
  explainability walkthroughs
\item
  multi-step assisted flows
\end{itemize}

This should compose:

\begin{itemize}
\tightlist
\item
  preference rows
\item
  alerts and notes
\item
  stepper or timeline language
\item
  explicit next actions
\end{itemize}

DL-005 proves the first canonical seam in this family through \texttt{e\allowbreak{}x\allowbreak{}p\allowbreak{}l\allowbreak{}a\allowbreak{}i\allowbreak{}n\allowbreak{}a\allowbreak{}b\allowbreak{}i\allowbreak{}l\allowbreak{}i\allowbreak{}t\allowbreak{}y\allowbreak{}(\allowbreak{})}, which turns the AI explainability doctrine into a structured guided-flow surface instead of leaving it as prose only. DL-008 promotes the reusable block seam through \texttt{g\allowbreak{}u\allowbreak{}i\allowbreak{}d\allowbreak{}e\allowbreak{}d\allowbreak{}F\allowbreak{}l\allowbreak{}o\allowbreak{}w\allowbreak{}(\allowbreak{})}, which standardizes calmer multi-step assistance, section rhythm, and next-action treatment without forcing every guided surface to masquerade as AI explainability.

\subsection{Inspector panel}\label{record-018-inspector-panel}

An inspector-style side panel is another strong candidate because it naturally combines:

\begin{itemize}
\tightlist
\item
  titled sections
\item
  supporting copy
\item
  current-value emphasis
\item
  compact but calm spacing
\end{itemize}

DL-006 now proves the first canonical seam in this family through \texttt{i\allowbreak{}n\allowbreak{}s\allowbreak{}p\allowbreak{}e\allowbreak{}c\allowbreak{}t\allowbreak{}o\allowbreak{}r\allowbreak{}(\allowbreak{})}, which turns inspector-panel rhythm into a reusable core component instead of leaving it trapped in a bespoke DOGFOOD detail box. DL-007 promotes the first block-level seam in this family through \texttt{i\allowbreak{}n\allowbreak{}s\allowbreak{}p\allowbreak{}e\allowbreak{}c\allowbreak{}t\allowbreak{}o\allowbreak{}r\allowbreak{}D\allowbreak{}r\allowbreak{}a\allowbreak{}w\allowbreak{}e\allowbreak{}r\allowbreak{}(\allowbreak{})}, which standardizes drawer-attached inspector composition in TUI shells instead of leaving existing inspector drawers to hand-compose strings.

\section{Block acceptance bar}\label{record-018-block-acceptance-bar}

A structure is ready to become a canonical block when:

\begin{itemize}
\tightlist
\item
  it recurs across more than one surface or product need
\item
  its interaction ownership is stable
\item
  its patterns are already canonical elsewhere
\item
  its lowering story is explicit
\item
  it improves reuse instead of trapping more quality inside one shell
\end{itemize}

\section{What not to do}\label{record-018-what-not-to-do}

\begin{itemize}
\tightlist
\item
  do not promote every nice-looking demo into a block
\item
  do not let blocks bypass canonical patterns
\item
  do not let blocks become storage, service, or app-domain assumptions
\item
  do not treat blocks as the first place where accessibility or localization is considered
\end{itemize}

\section{Relationship to the current cycle}\label{record-018-relationship-to-the-current-cycle}

DL-002 does not implement a blocks runtime.

This cycle exists to make the block layer explicit so future cycles can promote real recurring structures deliberately instead of accidentally.

